\section{Discrete random variables and their support}

A random variable is a function from the sample space to the real line. The next
question is: what kinds of values can this function produce? In this chapter we
focus on the discrete case, where the possible values form a finite or countable
set.

\begin{definitionbox}
A random variable \(X\) is \textbf{discrete} if the set of values it can take is
finite or countably infinite. Equivalently, its possible values can be written as
\[
x_1,x_2,x_3,\ldots
\]
with the list either ending after finitely many terms or continuing indefinitely.
\end{definitionbox}

The word discrete refers to the values of the random variable, not necessarily to
how complicated the original experiment is. The sample space may contain detailed
objects, but the random variable may record only a count, a sum, an indicator, or
a payoff.

\subsection*{The range and the support}

If \(X:\Omega\to\mathbb{R}\), then the set of values that \(X\) actually takes is
\[
X(\Omega)=\{X(\omega):\omega\in\Omega\}.
\]
This is the range of the function. In probability, we usually focus on values
that have positive probability.

\begin{definitionbox}
The \textbf{support} of a discrete random variable \(X\) is the set
\[
\operatorname{supp}(X)=\{x\in\mathbb{R}:P(X=x)>0\}.
\]
It is the set of values that occur with positive probability.
\end{definitionbox}

In elementary finite models, the range and support are often the same. But the
definition using positive probability is more robust, because it ignores values
that may be formally possible in notation but receive probability zero.

\subsection*{Example: number of heads}

Toss a coin \(n\) times and record the ordered sequence. The sample space is
\[
\Omega=\{H,T\}^n.
\]
Define
\[
X(\omega)=\text{the number of heads in }\omega.
\]
Then the possible values are
\[
0,1,2,\ldots,n.
\]
Thus
\[
\operatorname{supp}(X)=\{0,1,2,\ldots,n\},
\]
provided each sequence has positive probability.

This example shows a key point: \(\Omega\) contains sequences, but the support of
\(X\) contains numbers. The random variable turns a sequence into a count.

\subsection*{Example: sum of two dice}

Roll two dice and record the ordered pair:
\[
\Omega=\{1,2,3,4,5,6\}^2.
\]
Define the sum random variable
\[
S(i,j)=i+j.
\]
The smallest possible value is \(2\), and the largest possible value is \(12\).
Every integer between them can occur, so
\[
\operatorname{supp}(S)=\{2,3,4,5,6,7,8,9,10,11,12\}.
\]

Notice that the support has \(11\) values, while the original sample space has
\(36\) elementary outcomes. Many different pairs lead to the same sum. For
example,
\[
S^{-1}(7)=\{(1,6),(2,5),(3,4),(4,3),(5,2),(6,1)\}.
\]

\subsection*{Example: waiting time}

In the experiment of tossing a coin until the first head appears, define
\[
N=\text{the number of tosses needed to obtain the first head}.
\]
Then
\[
\operatorname{supp}(N)=\{1,2,3,\ldots\}.
\]
This support is countably infinite. The random variable is still discrete,
because its possible values can be listed.

\subsection*{Discrete does not mean finite}

A discrete random variable may have infinitely many possible values. The waiting
time \(N\) is the first important example. There is no largest possible waiting
time: success may occur on the first trial, the second trial, the third trial,
and so on. Still, the values are separated and countable:
\[
1,2,3,\ldots
\]
This is very different from a continuous variable such as a uniformly chosen
point in \([0,1]\), whose possible values cannot be listed in a sequence.

\begin{remarkbox}
Discrete means finite or countably infinite. Continuous means that probability is
spread over a continuum of values. The distinction concerns the structure of the
possible values of the random variable.
\end{remarkbox}

\subsection*{Events generated by values of a random variable}

Once \(X\) is defined, statements about \(X\) create events in the original
sample space. For example,
\[
\{X=3\}=\{\omega\in\Omega:X(\omega)=3\},
\]
\[
\{X\leq 3\}=\{\omega\in\Omega:X(\omega)\leq 3\}.
\]

For the number of heads in five tosses, the event \(\{X=3\}\) contains all
sequences with exactly three heads. For the sum of two dice, the event
\(\{S\leq 4\}\) contains all ordered pairs whose sum is \(2\), \(3\), or \(4\).

\subsection*{What the support tells us}

The support tells us where the probability law of the random variable lives. It
is the list of values for which we must assign probabilities. Once the support is
known, the next task is to determine
\[
P(X=x)
\]
for each \(x\) in the support.

\begin{theorembox}
To study a discrete random variable, first identify its support. The support is
the bridge between the original sample space and the numerical probability law of
\(X\).
\end{theorembox}
